\section{Introduction}\label{sec:introduction}

Let $X$ be a smooth connected complex projective variety of dimension
$2n+1$, where $n\geq1$. A \emph{contact structure} on $X$ is an exact
sequence of algebraic vector bundles
\begin{equation}\label{intro:contact-sequence}
 0\longrightarrow F\longrightarrow T_X
 \xrightarrow{\ \theta\ }L\longrightarrow0,
\end{equation}
with $L$ a line bundle, such that the Levi pairing
\[
 F\otimes F\longrightarrow L,
 \qquad (v,w)\longmapsto\theta([v,w])
\]
is everywhere nondegenerate. In a local frame of $L$, write $\alpha$
for the scalar one-form representing $\theta$. The condition is
equivalent to $\alpha\wedge(d\alpha)^n$ being nowhere zero. We call
$(X,F)$ a \emph{contact Fano manifold} if $-K_X$ is ample.

For a simple complex Lie algebra $\mathfrak g$, let
$\mathcal O_{\min}(\mathfrak g)$ be its nonzero nilpotent adjoint orbit
of smallest dimension, and put
\[
 Y_{\mathfrak g}=\PP(\mathcal O_{\min}(\mathfrak g)).
\]
This is the unique closed adjoint orbit in $\PP(\mathfrak g)$. Its
canonical contact structure is obtained by projectivizing the
Kirillov--Kostant--Souriau symplectic cone.

\begin{theorem}\label{thm:main}
Let $(X,F)$ be a smooth connected complex projective contact Fano
manifold of dimension $2n+1$, with $n\geq1$, and let $L=T_X/F$.
There exist a finite-dimensional simple complex Lie algebra
$\mathfrak g$ and an algebraic isomorphism
\[
 \varphi:X\xrightarrow{\ \sim\ }Y_{\mathfrak g}
\]
whose differential carries the given distribution $F$ to the canonical
contact distribution on $Y_{\mathfrak g}$. Moreover,
\[
 L\simeq\varphi^*\bigl(\OO_{\PP(\mathfrak g)}(1)|_{Y_{\mathfrak g}}\bigr).
\]
\end{theorem}

Throughout the paper, $\PP(E)$ parametrizes lines in $E$, with
tautological subbundle $\OO_{\PP(E)}(-1)$. The two models separated
by the Picard reduction are
\[
 \begin{array}{c|c|c}
 X&L&\mathfrak g\\ \hline
 \PP^{2n+1}&\OO(2)&\mathfrak{sp}_{2n+2}(\CC)\\[2pt]
 \PP(T^*\PP^{n+1})&\OO(1,1)|_X&\mathfrak{sl}_{n+2}(\CC).
 \end{array}
\]
In the second row, $X$ is the incidence divisor of points and
hyperplanes in $\PP^{n+1}\times(\PP^{n+1})^*$. In the first, the
adjoint embedding is the quadratic Veronese embedding, so its
hyperplane bundle is $\OO(2)$. The proof then treats the remaining
case $\Pic(X)=\ZZ[L]$.

Two consequences will be proved after the classification in
Section~\ref{hom:section}.

\begin{corollary}[K\"ahler--Einstein existence]\label{thm:contact-fano-ke}
Let $(X,F)$ be a smooth connected complex projective contact Fano
manifold of complex dimension $2n+1$, where $n\geq1$. Then $X$ admits a
K\"ahler--Einstein metric of positive scalar curvature for its given
complex structure.
\end{corollary}

For projective contact manifolds without a Fano hypothesis, the
results of Demailly and Kebekus--Peternell--Sommese--Wi\'sniewski
provide a reduction to which Theorem~\ref{thm:main} applies.

\begin{corollary}[Projective contact manifolds]\label{thm:projective-contact}
Let $(X,F)$ be a smooth connected complex projective contact manifold
of dimension $2n+1$, where $n\geq1$. Its second Betti number determines
the following alternatives.
\begin{enumerate}
\item If $b_2(X)=1$, then $(X,F)$ is contact-isomorphic to the adjoint
variety $Y_{\mathfrak g}$, with its canonical contact distribution, for
a simple complex Lie algebra $\mathfrak g$.
\item If $b_2(X)\geq2$, then there is a smooth connected complex
projective manifold $Z$ of dimension $n+1$ and a biholomorphism
\[
 X\simeq\PP(T^*Z).
\]
Here $\PP(T^*Z)$ parametrizes lines in cotangent spaces, equivalently
hyperplanes in tangent spaces. This assertion identifies the underlying
manifolds, without imposing compatibility with the given distribution $F$.
\end{enumerate}
Conversely, every adjoint variety of complex dimension $2n+1$ and every
projectivized cotangent bundle in the second alternative admits a contact
structure. Thus these are precisely the underlying manifolds of smooth
connected complex projective contact manifolds in these dimensions.
\end{corollary}

The Riemannian application concerns a smooth manifold of real
dimension $4m$, $m\geq2$, whose metric has holonomy contained in
$\Sp(m)\Sp(1)$ in its standard representation. Such a metric is
called \emph{quaternionic-K\"ahler}; \emph{positive} here means that
its scalar curvature is positive. The compact symmetric positive
quaternionic-K\"ahler spaces are the \emph{Wolf spaces}.

\begin{corollary}\label{thm:quaternionic}
Let $(M^{4m},g)$ be a closed connected smooth Riemannian manifold,
where $m\geq2$, with holonomy contained in $\Sp(m)\Sp(1)$ in the
standard representation and with positive scalar curvature. Then
\[
 \nabla\Rm_g=0.
\]
Moreover, $(M,g)$ is homothetic to a compact symmetric positive
quaternionic-K\"ahler Wolf space.
\end{corollary}

Theorem~\ref{thm:main} resolves the smooth contact-Fano homogeneity
conjecture positively. Corollary~\ref{thm:quaternionic} gives the
positive resolution of the Riemannian LeBrun--Salamon conjecture in
every real dimension $4m$ with $m\geq2$. The link is the positive
twistor construction: the twistor space is contact Fano, so the
general algebraic theorem applies to it. Simple connectivity of the
contact Fano twistor space and its $S^2$-bundle structure show that
the original manifold is simply connected. The metric is recovered
from the resulting complex twistor space in Section~\ref{tw:section}.

\subsection*{History and prior work}

The conjectural classification has two related settings.  Wolf's work
relates compact quaternionic symmetric spaces to homogeneous complex contact
spaces~\cite{Wolf1965}; the adjoint-variety description of their twistor
fibration is recalled in~\cite[Section~3.9, pp.~105--106]{LandsbergManivel2004}.
The twistor construction of Salamon, independently due to B\'erard-Bergery,
associates a contact Fano manifold to a closed connected positive
quaternionic-K\"ahler manifold~\cite{Salamon1982,LeBrun1995}.  LeBrun's converse characterizes the
contact Fano manifolds arising in this way by the existence of a
K\"ahler--Einstein metric~\cite[Theorem~A]{LeBrun1995}.  The finiteness and
rigidity program of LeBrun and Salamon~\cite{LeBrun1993,LeBrunSalamon1994}
provides the historical setting for the homogeneity conjecture.  The general
algebraic formulation, stated for example in
\cite[Conjecture~1.1]{Buczynski2010}, concerns every contact Fano manifold.

The projective-contact reduction has two ingredients. Kebekus,
Peternell, Sommese and Wi\'sniewski classify projective contact
manifolds with non-nef canonical bundle as either Fano with $b_2=1$
or projectivized cotangent bundles \cite[Theorem~1.1]{KPSW2000}.
Demailly's integrability theorem excludes pseudoeffectivity of the
canonical bundle of a compact K\"ahler contact manifold; his
Corollaries~3--4 give the reduction used in
Corollary~\ref{thm:projective-contact}
\cite[Corollaries~2--4]{Demailly2000}.
Within the Fano class, Mori's ample-tangent theorem and the
Kobayashi--Ochiai index theorem identify the incidence and
projective-space alternatives \cite{Mori1979,KO1973}. In the
remaining primitive case, Kebekus establishes the covering contact
lines and the spanning of their tangent directions
\cite{Kebekus2001}. For $n\geq2$, the corrected contact-line theorem gives a smooth,
possibly disconnected, Legendrian variety of directions
\cite[corrected version, Theorem~3.1 and Remark~3.2]{Buczynski2010}.
The case $n=1$ is treated directly in Section~\ref{lines:section}.
These results identify the geometric objects from which the
smoothing, tensor descent and diagonal arguments below begin.


Earlier low-dimensional classifications were obtained by Poon and Salamon
for compact positive quaternionic-K\"ahler manifolds in real dimension
eight~\cite{PoonSalamon1991}, and by Ye and Druel for contact Fano manifolds in
complex dimensions three and five, respectively
\cite{Ye1994,Druel1998}.
Several classification results impose additional geometric or
group-theoretic hypotheses.  Beauville's contact moment-map approach gives adjoint
classification when the contact linear-system map is generically finite and
the contact automorphism group is reductive~\cite[Theorem~0.1]{Beauville1998}.
Yasukura proves adjoint classification for compact connected complex
contact manifolds whose contact line bundle defines a holomorphic
immersion~\cite[Theorem~1.1]{Yasukura2023v12}. The first-jet generation
obtained below supplies such an immersion; we give the resulting
contact moment-map classification directly.
Buczy\'nski and Wi\'sniewski classify the closed positive quaternionic-K\"ahler
manifolds in real dimensions twelve and sixteen, and the contact Fano
manifolds in complex dimensions three through nine with reductive automorphism
group~\cite[Theorems~1.1--1.2]{BW2022}.
Occhetta, Romano, Sol\'a Conde, and Wi\'sniewski prove the adjoint
classification in the primitive Picard case under reductivity and the rank
bound $r\geq\max\{2,(n-3)/2\}$ in complex dimension $2n+1$,
where $r$ is the rank of the contact automorphism group
\cite[Theorem~6.1]{OcchettaEtAl2021}.
On the Riemannian side, Brendle and Semmelmann establish symmetry for
compact positive quaternionic-K\"ahler manifolds of real dimension at least
eight under nonnegative sectional curvature
\cite[Corollary~2]{BrendleSemmelmann2025}.

The use of contact lines and two-line loci has substantial antecedents.
Buczy\'nski attributes the study of loci swept out by two contact lines to
Wi\'sniewski~\cite[Introduction]{Buczynski2010} and develops their divisors, duality, and
associated bilinear forms~\cite[Introduction and Theorem~1.3]{Buczynski2010}.
The Legendrian secant geometry of Landsberg and Manivel
\cite[Proposition~17(2) and its proof]{LandsbergManivel2007}
is used in Buczy\'nski's Proposition~3.5, which also records the
disjoint-join statement~\cite{Buczynski2010}.
Here a suitable pair of contact lines is smoothed, two vector lifts are
extended across the unmarked degree-two boundary, and their common contact
value descends to a symmetric tensor on the square.  Its diagonal jets
generate ordinary first jets of the contact line bundle at every point,
which yields transitivity of contact automorphisms.  The final jet criterion
also applies to compact complex contact manifolds satisfying its tensor
hypothesis, without a Fano assumption.

Kobayashi's preprint claims the Wolf-space classification for complete
locally irreducible positive quaternionic-K\"ahler manifolds of real
dimension at least eight
\cite[Theorem~5.3]{Kobayashi2008} and states K\"ahler--Einstein existence
for arbitrary contact Fano manifolds as a separate conjecture
\cite[Conjecture~1.3]{Kobayashi2008}. Corollary~\ref{thm:contact-fano-ke}
establishes this existence for the manifolds in Theorem~\ref{thm:main}.
Earlier versions of Yasukura's preprint asserted unrestricted contact-Fano
and quaternionic-K\"ahler classifications
\cite[Theorem~I and Corollary~II]{Yasukura2009v3},
\cite[Theorems~C--D and Corollary~E]{Yasukura2022v11}.
The author's later comments identify an error in Theorem~3.2(ii) of the
former version~\cite{Yasukura2009v4}, and identify Theorem~A and
Corollary~B of the latter version as false~\cite{Yasukura2023v12}.


\subsection*{The mechanism of the proof}

The \emph{contact degree} of a map $f:\PP^1\to X$ is
$\deg f^*L$. We call such a map a \emph{contact line} if its contact
degree is one and a \emph{conic} if its contact degree is two; these
terms impose no embedding assumption.

The objective is to produce enough global vector fields preserving
$F$ to span every tangent space. After the two explicit Picard
alternatives have been handled, we work with $\Pic(X)=\ZZ[L]$ and
construct a symmetric section
\begin{equation}\label{intro:global-section}
 B\in H^0(X\times X,L\boxtimes L),\qquad
 L\boxtimes L=\operatorname{pr}_1^*L\otimes\operatorname{pr}_2^*L,
\end{equation}
where $\operatorname{pr}_1,\operatorname{pr}_2$ are the projections
from $X\times X$. Near two coincident points the section vanishes
to second order, with leading term given by the square of the
contact form. This local behavior will force the global contact
fields to span. The construction has four stages.

First, at a general point choose two free contact lines, meaning
that their pulled-back tangent bundles are globally generated.
Their positive-degree tangent-bundle summands span $F_x$ and their tangent
vectors have nonzero Levi pairing. Their union can be smoothed
while its two outer values move. The key calculation concerns
the contrary, fixed-end problem: a first-order smoothing keeping
both values fixed would force the Levi pairing to vanish. The
fixed-end obstruction space is one-dimensional, so varying the
smoothing parameter supplies exactly the missing direction for
two-point evaluation. The resulting conic $g:\PP^1\to X$ is
\emph{very free}, meaning that every summand of $g^*T_X$ has
positive degree (Proposition~\ref{sm:very-free}). We keep all
nearby parametrized maps, not just a chosen one-parameter
smoothing, because the later diagonal calculation varies both
their values and their velocities.

Second, compactify these conics in the unmarked stable-map
space. Its universal-curve square parametrizes a source curve
together with two variable evaluation points. On a smooth source,
the coordinate-independent tensor
\[
 (z-w)^2\partial_z\otimes\partial_w
\]
gives two vector-valued sections, one retaining the tangent vector
in the first factor and one retaining it in the second. Applying
the contact form to both factors gives their common contact value.
The singular unmarked sources have exactly two degree-one
components. At a node of thickness $r$, the contact differential
vanishes to order at least $r$, canceling each possible pole of
the tensor. Both vector-valued sections consequently extend over
the proper family (Proposition~\ref{ext:moving-pair}).

Third, the common contact value descends from this proper family
to the section $B$ in \eqref{intro:global-section}. The two vector
lifts test contact lines in the two factors of $X\times X$.
On every proper integral component dominating a test line, the
relevant lift has constant projection to the trivial quotient of
the pulled-back tangent bundle. This forces every generic sheet
of the finite contact-value image to split over that line. A test
line meeting a hypothetical branch divisor transversely would
instead produce a sheet of degree greater than one. Purity and
simple connectedness then complete the descent
(Proposition~\ref{des:descent}).

Finally, on the full neighborhood of $g$ the section satisfies
\[
 B(h(z),h(w))=(z-w)^2\,
 \theta(h'(z))\otimes\theta(h'(w)).
\]
The available value and velocity variations make its constant
and linear diagonal terms vanish and identify the quadratic
symbol with $\theta^2$. Each of these intrinsic identities
extends from a nonempty open subset to all of $X$. Symmetry then
determines the cubic Taylor term in local coordinates and a line
frame. The quadratic and cubic terms imply that global sections
of $L$ generate first jets at every point. Contact Hamiltonians
therefore span every tangent space, giving contact transitivity
(Theorem~\ref{jet:criterion}). The contact moment map identifies
the resulting homogeneous manifold, with its given distribution,
with a simple adjoint variety (Proposition~\ref{hom:adjoint}).

Two parts of this mechanism have independent formulations:
contact-value descent from a proper family with the two vector
lifts, and the symmetric-jet criterion on a connected compact
complex contact manifold. The latter does not require projectivity
or Fano positivity. Global diagonal identities are what pass
from the initially general points to every point. After
transitivity is proved, Corollary~\ref{jet:good-curves} supplies
through every point in the primitive Picard case a conic transverse
to $F$ with $u^*F\simeq\OO(1)^{\oplus2n}$.

\subsection*{Conventions and organization}

All schemes, algebraic spaces and stacks are over $\CC$.
Curve families have their
scheme-theoretic fibers. ``General'' means outside a proper
Zariski-closed subset unless an analytic open set is explicitly
specified. For bundles $E_1,E_2$ on $X$, we similarly write
$E_1\boxtimes E_2=\operatorname{pr}_1^*E_1\otimes
\operatorname{pr}_2^*E_2$.

Section~\ref{red:section} proves the Picard reduction.
Section~\ref{lines:section} selects the pair of contact lines,
and Section~\ref{sm:section} smooths it to a very free conic.
Section~\ref{ext:section} extends the two moving-point tensors
over the unmarked boundary; Section~\ref{des:section} descends
their contact value. Section~\ref{jet:section} turns its diagonal
jets into contact transitivity, and Section~\ref{hom:section}
identifies the adjoint model and completes Theorem~\ref{thm:main}.
Section~\ref{tw:section} gives the twistor construction and
normalized metric recovery proving Corollary~\ref{thm:quaternionic}.
